\begin{abstract}
Transient hardware faults threaten large language model (LLM) inference, where autoregressive decoding and multi-GPU execution allow errors to propagate across tokens and devices. Existing model-level fault-injection tools provide limited support for distributed and quantized inference, often overlook computation-dependent fault exposure and communication faults, and incur repeated host--device transfers for bit mutation. We present vLLM-FI, a fault-injection and reliability-evaluation framework integrated with vLLM. The framework models memory, computation, and communication faults, accounting for computational exposure beyond output tensor size and perturbing software-visible communication buffers. Configuration-driven targeting coordinates injection across distributed workers, while customized CUDA operators flip bits in place on GPU-resident tensors. Support for tensor and pipeline parallelism, quantized representations, and standardized downstream-task evaluation enables systematic reliability studies under practical inference configurations. In calibration experiments, vLLM-FI matches MRFI's output-level fault effects with a maximum absolute difference of 0.62 percentage points in fault-injection correct rate. Across the evaluated all-layer projection-injection workloads, vLLM-FI achieves a $1.81\times$--$2.41\times$ end-to-end speedup over MRFI. With the inference runtime held fixed, GPU-resident injection delivers up to a $5.22\times$ speedup in total experiment runtime over an MRFI-style CPU-mediated backend under compute-fault injection. These results demonstrate lower experimental cost while preserving agreement with the reference injector under the tested settings, enabling broader reliability evaluation across fault models, tasks, quantization methods, and parallel configurations.
\end{abstract}

\begin{IEEEkeywords}
Fault injection, reliability evaluation, large language models, distributed inference, GPU computing, quantization.
\end{IEEEkeywords}
